%% Centerpiece chapter, adapted from p01-loop-extrapolation
\subsection{The Loop-Depth Sweep}
\label{sec:loopsweep}

\subsubsection{Method}

\paragraph{Evaluation protocol.}  After training, we freeze the weights
and evaluate with loop counts $n \in \{1, 2, 4, 8, 12, 16\}$ on 256
fresh TinyStories sequences of length 256. For each $n$ we record mean
cross-entropy loss and perplexity, and generate continuations from a
fixed prompt for qualitative comparison. Sampling uncertainty is
quantified with bootstrap 95\% confidence intervals (10,000 percentile
resamples over the 256 sequences); all paired comparisons below use
per-sequence paired resampling. The $n=4$ setting matches
training; all others are pure test-time extrapolation or reduction.

\subsubsection{Results}

Table~\ref{tab:sweep} reports the sweep. Three regimes are visible.

\begin{table*}[t]
\centering
{\footnotesize\setlength{\tabcolsep}{6pt}
\begin{tabular}{lcccccc}
\toprule
Loops $n$ & 1 & 2 & 4 (trained) & 8 & 12 & 16 \\
\midrule
Loss & 3.380 & 2.414 & \textbf{1.484} & 1.590 & 1.591 & 1.591 \\
\phantom{Loss} 95\% CI & [3.343, 3.419] & [2.384, 2.446] & \textbf{[1.454, 1.516]} & [1.559, 1.622] & [1.560, 1.624] & [1.560, 1.623] \\
Perplexity & 29.38 & 11.18 & \textbf{4.41} & 4.90 & 4.91 & 4.91 \\
\phantom{Perplexity} 95\% CI & [28.30, 30.52] & [10.84, 11.54] & \textbf{[4.28, 4.55]} & [4.75, 5.06] & [4.76, 5.07] & [4.76, 5.07] \\
\bottomrule
\end{tabular}}
\caption{Test-time loop sweep: mean loss and perplexity on 256 held-out
TinyStories sequences (length 256). 95\% confidence intervals are
bootstrap percentile intervals (10,000 resamples over the 256 sequences).
Bold marks the training depth.}
\label{tab:sweep}
\end{table*}

\paragraph{Recurrence helps up to the trained depth.} Cutting loops
below training is catastrophic for quality: a single loop yields
perplexity 29.38, and doubling to 2 loops more than halves it to 11.18.
The trained depth of 4 reaches 4.41. The recurrent iterations are doing
substantial, non-redundant work.

\paragraph{Extrapolation yields no gain.} Running 8 loops---twice the
training count---\emph{hurts} slightly (4.41 $\rightarrow$ 4.90), and
further loops to 12 and 16 change nothing (4.91, 4.91). The dynamics do
not continue to refine; they appear to converge to a fixed point near,
but slightly worse than, the trained-depth output. The degradation is
not sampling noise: the 95\% bootstrap intervals on the two means do not
overlap ([4.28, 4.55] at 4 loops vs.\ [4.75, 5.06] at 8 loops), and the
paired per-sequence difference is $+0.106$ loss with interval
$[+0.101, +0.110]$. The 8-to-12 difference, by contrast, is $+0.001$
loss---negligible in magnitude---confirming a genuine
plateau rather than continued drift.

\paragraph{Generations stay coherent.} Qualitatively, samples at all
loop counts are grammatical story continuations; over-looping does not
cause collapse or repetition. The degradation is in token-level
likelihood, not in catastrophic failure modes.

\subsubsection{Discussion}

Our central finding is that the recurrent dynamics are \emph{calibrated}
to the training iteration count. This is consistent with the model
learning an approximate fixed-point iteration whose contraction
properties hold for the trained horizon: within it, each loop refines;
beyond it, the iteration has converged and additional steps add only
noise.

This cautions against treating loop count as a free test-time scaling
knob. If extrapolation is desired, it likely requires training-time
support---e.g., randomized loop counts per batch, or an explicit
halting objective---rather than emerging by default. Notably, our
results do not contradict the broader test-time scaling literature:
sampling-based methods \citep{snell2024scaling} scale a different axis
(number of attempts) than recurrent refinement (depth of a single
attempt).

A limitation is scale: our study uses one 44M model on one dataset.
Within that scope, sampling uncertainty is now quantified---the key
4-to-8-loop degradation survives a paired bootstrap test---but
across-run uncertainty (different seeds, different data orderings)
remains unmeasured and is the natural next experiment.
Whether larger recurrent models, or training with loop-count
randomization, recover extrapolation is left to future work. A second
limitation is metric scope: perplexity may miss qualitative gains from
extra loops on harder reasoning tasks, which TinyStories does not
exercise.

\subsubsection{Conclusion}

On a 44M recurrent transformer trained at 4 loops, test-time loop
extrapolation does not deliver free gains: quality saturates at the
training depth and slightly degrades beyond it (4.41 $\rightarrow$
4.90 perplexity at 8 loops) before plateauing.
Recurrence is powerful within its trained horizon and yields nothing
further outside it.
We release our evaluation protocol to enable comparable sweeps at
larger scales.
